% Source-derived chapter 02: Representation-theoretic preliminaries
% Original source: canonical-representations-surface-groups
\chapter{Representation-theoretic preliminaries}
\label{canonical-representations-surface-groups-chapter-02}
\source{canonical-representations-surface-groups}

\begin{remark}[Source section]
\label{canonical-representations-surface-groups-chapter-02-source}
\source{canonical-representations-surface-groups}
\source{canonical-representations-surface-groups}
This chapter reproduces the corresponding section of the retrieved
LaTeX source, including its definitions, statements, examples, and
proofs. It has not yet been translated into Lean.
\notready
\end{remark}

% The source section title was: Representation-theoretic preliminaries
\label{section:representation-preliminaries}
\source{canonical-representations-surface-groups}

In this section, we give group-theoretic constructions which we will use to
analyze representations with finite mapping class group orbit.
In \autoref{subsection:basic-properties},
we verify basic properties of MCG-finite representations. 
%We collect facts about deformations of representations in
%\autoref{subsection:tangent-space}.
In \autoref{subsection:mcg-finite-construction}, given an irreducible representation
$\rho: \pi_1(\Sigma_{g,n},x) \to \on{GL}_r(\mathbb C)$ whose conjugacy class is fixed by a finite-index subgroup $\Gamma\subset\on{PMod}_{g,n+1}$, we construct a representation
$\widetilde{\rho}:  \Gamma \to \on{PGL}_r(\mathbb C)$,
which will be analyzed throughout this paper.
In \autoref{subsection:lifting-projective},
we show how to lift certain projective representations of the fundamental groups of families of curves to honest representations into
$\on{GL}_r(\mathbb C)$,
after passing to a suitable cover. Finally, in
\autoref{subsection:unitary-on-fibers}, we prove some structural results about the representations we've constructed.

We will use the following lemma throughout, to connect properties of mapping class groups to geometry.
\begin{lemma}
\source{canonical-representations-surface-groups}
\notready
	\label{lemma:mgn-fundamental-group}
\source{canonical-representations-surface-groups}
	For $m \in \mathscr M_{g,n}$ a basepoint and $m'\in \mathscr M_{g, n+1}$ a lift of $m$, there are isomorphisms $\pi_1(\mathscr M_{g,n}, m) \simeq \on{PMod}_{g,n}$, $\pi_1(\mathscr M_{g,n+1}, m') \simeq \on{PMod}_{g,n}$ such that the diagram 
	$$\xymatrix{
	1 \ar[r] & \pi_1(\Sigma_{g,n}) \ar[r] \ar@{=}[d]& \pi_1(\mathscr{M}_{g,n+1}, m') \ar[r] \ar[d]^\sim & \pi_1(\mathscr{M}_{g,n}, m) \ar[r] \ar[d]^\sim & 1 \\
	1 \ar[r] & \pi_1(\Sigma_{g,n}) \ar[r] & \on{PMod}_{g,n+1} \ar[r] & \on{PMod}_{g,n}\ar[r] & 1
	}$$
	commutes, where the vertical maps are given by these isomorphisms, the top row is the exact sequence of homotopy groups induced by the forgetful map $\mathscr{M}_{g,n+1}\to \mathscr{M}_{g,n}$,  and the bottom row is the Birman exact sequence \cite[p.~98]{farbM:a-primer}. 
\end{lemma}
\begin{proof}
This follows from the contractibility of the universal cover of $\mathscr{M}_{g,
n}$, and the fact that $\mathscr{M}_{g,n}$ is the quotient of its universal
cover by the properly discontinuous action of $\on{PMod}_{g,n}$. See
\cite[\S10.6.3 and p. 353]{farbM:a-primer}. There is a choice involved here, namely: if $C$ is the Riemann surface associated to the point $m\in \mathscr{M}_{g,n}$, one must choose a homeomorphism between $C$ and our reference surface $\Sigma_{g,n}$, and similarly with $m'$. Changing this homeomorphism replaces the vertical isomorphisms by conjugate isomorphisms.
%The case $n = 0$ is explained in \cite[p. 361]{farbM:a-primer}, where
%	the real content is that the universal cover of $\mathscr M_{g}$ is
%	contractible, as shown in \cite[Theorem 10.6]{farbM:a-primer}, and that
%	the action of $\on{Mod}_g$ on this universal cover is properly
%	discontinuous \cite[Theorem 12.2]{farbM:a-primer}.
%	In the general case $n \geq 0$, we can use \cite[\S 10.6.3]{farbM:a-primer}.
%	in place of \cite[Theorem 10.6]{farbM:a-primer} which shows the
%	universal cover of $\mathscr M_{g,n}$ is contractible
%	as well as \cite[p. 353]{farbM:a-primer} to generalize the proper
%	discontinuity of \cite[Theorem
%	12.2]{farbM:a-primer}, to the case of punctured surfaces.
\end{proof}
The subgroup $ \pi_1(\Sigma_{g,n})\subset \on{PMod}_{g,n+1} $ is often referred to as the \emph{point-pushing subgroup}.

\subsection{Basic properties of MCG-finite representations}\label{subsection:basic-properties}
\source{canonical-representations-surface-groups}

To acquaint the reader with MCG-finite representation, we now spell out some of
their basic properties.

\begin{proposition}
\source{canonical-representations-surface-groups}
\notready\label{proposition:basic-properties}
\source{canonical-representations-surface-groups}
	The direct sum and tensor product of two MCG-finite representations of $\pi_1(\Sigma_{g,n})$ is MCG-finite. Any semisimple subquotient of an MCG-finite representation is MCG-finite.
\end{proposition}
\begin{proof}
The first statement, about direct sums and tensor products, is clear. The second is immediate from \cite[Lemma 2.2.1]{lawrence2019representations}.
%Letting $\Gamma\subset \on{Mod}_{g,n}$ be the (finite-index) stabilizer of the conjugacy class of $\rho$, observe that $\Gamma$ acts on the set of irreducible subquotients of $\rho$. But this set is finite; hence the stabilizer of $\sigma$ is finite index in $\Gamma$, and hence in $\on{Mod}_{g,n}$.
\end{proof}

We next show that restrictions of representations of $\on{Mod}_{g,n+1}$ to the
point-pushing subgroup are MCG-finite. See also \cite[Lemma
2.2]{kasahara2015visualization} for a related result with analogous proof. 

\begin{proposition}
\source{canonical-representations-surface-groups}
\notready\label{proposition:MCG-finite-index-rep}
\source{canonical-representations-surface-groups}
	Let $\Gamma\subset \on{Mod}_{g,n+1}$ be a finite index subgroup
	containing the point-pushing subgroup $\pi_1(\Sigma_{g,n})\subset
	\on{Mod}_{g,n+1}$. Let $$\rho:\Gamma\to \on{GL}_r(\mathbb{C})$$ be a representation. Then $\rho|_{\pi_1(\Sigma_{g,n})}$ is MCG-finite.
\end{proposition}
\begin{proof}
	By replacing $\Gamma$ with 
	$\Gamma \cap \on{PMod}_{g,n+1},$
	we may assume $\Gamma \subset \on{PMod}_{g,n+1}$,
	the pure mapping class group.  We claim that the image of $\Gamma$ in
	$\on{PMod}_{g,n}$ (under the map $\on{PMod}_{g,n+1}\to \on{PMod}_{g,n}$
	arising from the Birman exact sequence) stabilizes the conjugacy class
	of $\rho$.
	It suffices to show that for each $\gamma\in \Gamma$, 
\begin{align*}
	\rho^\gamma:  \pi_1(\Sigma_{g,n})& \rightarrow \on{GL}_r(\mathbb C) \\
	g & \mapsto \rho(\gamma g \gamma^{-1})
\end{align*}
is conjugate to $\rho$. Indeed, $\rho^\gamma(h) = \rho(\gamma) \rho(h)
\rho(\gamma)^{-1}$.
\end{proof}

We later give a partial converse to \autoref{proposition:MCG-finite-index-rep},
for irreducible representations, in \autoref{corollary:converse-to-MCG-rep}.
We now give a geometric counterpart to
\autoref{proposition:MCG-finite-index-rep},
which is closely related to \cite[Theorem A1]{CH:isomonodromic}.
\begin{proposition}
\source{canonical-representations-surface-groups}
\notready
\label{proposition:MCG-finite-family-of-curves}
\source{canonical-representations-surface-groups}
%	Let $\mathscr{C}\to \mathscr{M}$ be a versal family of $n$-pointed
%	curves of genus $g$, as in \autoref{notation:versal-family}. 
Let $\pi^\circ: \mathscr{C}^\circ\to \mathscr{M}$ be a punctured versal family
of $n$-pointed genus $g$ curves, as in \autoref{notation:versal-family}. In
particular, $\mathscr M \to \mathscr M_{g,n}$ is dominant and \'etale.
	Let $C^\circ$ be a fiber of $\pi^\circ$. If
	$$\rho:\pi_1(\mathscr{C}^\circ)\to \on{GL}_r(\mathbb{C})$$ is a representation, $\rho|_{\pi_1(C^\circ)}$ is MCG-finite.
\end{proposition}
We prove \autoref{proposition:MCG-finite-family-of-curves}, in
\autoref{subsubsection:proof-family-of-curves},
but we first require some well-known lemmas about families of curves. 
\begin{lemma}
\source{canonical-representations-surface-groups}
\notready
	\label{lemma:finite-index-from-dominant-map}
\source{canonical-representations-surface-groups}
	For $\mathscr M$ a scheme and any dominant \'etale map
	$\mathscr M \to \mathscr M_{g,n}$,
	$\on{im}(\pi_1(\mathscr M) \to \pi_1(\mathscr M_{g,n}))$
	has finite index in 
$\pi_1(\mathscr M_{g,n})\simeq \on{PMod}_{g,n}$.
\end{lemma}
\begin{proof}
	After pulling back to a finite \'etale cover $\mathscr
	N \to \mathscr M_{g,n}$ which is a scheme (such exists by e.g.~\cite[Proposition 2.3.4]{pikaart-dejong}),
	we may assume there is a dominant \'etale map of schemes
	$\alpha: \mathscr M \to \mathscr N$. It is enough to show $\on{im}(\pi_1(\mathscr M) \to \pi_1(\mathscr N))$
	has finite index in $\pi_1(\mathscr N)$, which follows from \cite[Lemma
	4.19]{debarre:higher-dimensional}.
	The basic idea here is to pass to opens over which the map is finite
	flat, and use that passing to opens will preserve the property that the
	image of the map of fundamental groups has finite index.
%By generic finiteness and upper semicontinuity of fiber rank, there exists a
%dense open subscheme $\mathscr N^\circ \subset \mathscr N$ over which
%$\alpha: \mathscr M \to \mathscr N$
%is finite flat.
%Let $\mathscr M^\circ := \alpha^{-1}(\mathscr N^\circ)$.
%Since $\mathscr M^\circ \to \mathscr N^\circ$ is finite \'etale,
%$\pi_1(\mathscr M^\circ) \to \pi_1(\mathscr N^\circ)$ has finite index.
%
%Therefore, 
%since $\pi_1(\mathscr M^\circ) \to \pi_1(\mathscr N^\circ)$ has finite index,
%in order to show 
%$\pi_1(\mathscr M) \to \pi_1(\mathscr N)$ has finite index, it is 
%enough to show
%$\pi_1(\mathscr M^\circ) \to \pi_1(\mathscr M)$ and
%$\pi_1(\mathscr N^\circ) \to \pi_1(\mathscr N)$ are surjective.
%This holds by the following fact:
%If $X$ is a smooth connected scheme, any dense open subscheme $U \subset X$ 
%has complement of real codimension $2$, so $\pi_1(U) \to \pi_1(X)$ is
%surjective, as any loop in
%$X$ is homotopic to one in $U$. 
\end{proof}

%We next describe an exact sequence on fundamental groups associated to the
%fibration $\mathscr C^\circ \to \mathscr M$. 
%we will use the associated Hochschild-Serre spectral sequence.
\begin{lemma}
\source{canonical-representations-surface-groups}
\notready\label{lemma:SES-pi1-versal-curves}
\source{canonical-representations-surface-groups}
If $\mathscr{C} \to \mathscr{M}$ is a versal family of $n$-pointed curves of
genus $g$, with $3g - 3 + n \geq 0$, as in \autoref{notation:versal-family},
so that $C^\circ$ is a fiber of $\pi^\circ$,
the sequence
\begin{equation}
	\label{equation:fibration-sequence}
\source{canonical-representations-surface-groups}
	\begin{tikzcd}
		1 \ar {r} & \pi_1({C}^\circ)  \ar {r} &
		\pi_1(\mathscr{C}^\circ) \ar {r} & \pi_1(\mathscr{M}) \ar {r} &
		1
\end{tikzcd}\end{equation}
is exact.
\end{lemma}
\begin{proof}
Except for injectivity, the result is immediate from the long exact sequence in
homotopy groups and the fact that ${C}^\circ$ is connected. Injectivity holds in
the universal case where $\mathscr{M}=\mathscr{M}_{g,n}$ is the Deligne-Mumford
moduli stack of curves and $\mathscr{C}^\circ = \mathscr{C}^\circ_{g,n}$ is the
universal punctured curve, as the universal cover of $\mathscr{M}_{g,n}$ is contractible. 
We next show the map $\pi_1(C^\circ) \to
\pi_1(\mathscr C^\circ)$ is injective in the general case via pullback.
The composite map $C^\circ \to \mathscr C \to \mathscr{C}^\circ_{g,n}$
is the natural inclusion of $C^\circ$ as a fiber of $\mathscr{C}^\circ_{g,n}$
over $\mathscr M_{g,n}$.
Hence, the composition $\pi_1(C^\circ) \to \pi_1(\mathscr C^\circ) \to \pi_1(\mathscr
C^\circ_{g,n})$ is injective, and so the first map is injective.
\end{proof}


\subsubsection{}
\label{subsubsection:proof-family-of-curves}
\source{canonical-representations-surface-groups}
\begin{proof}[Proof of {\autoref{proposition:MCG-finite-family-of-curves}}]
	The exact sequence \eqref{equation:fibration-sequence}
induces an outer action of $\pi_1(\mathscr M)$ on
$\pi_1(C^\circ)$, which factors through the outer action of $\on{PMod}_{g,n}$ on
$\pi_1(C^\circ)$ induced by the Birman exact sequence.
Specifically, this outer action sends $\gamma \in \on{PMod}_{g,n}$
to the action of conjugation by $\widetilde{\gamma}$ on $\pi_1(C^\circ)$, for
$\widetilde{\gamma}\in \on{PMod}_{g,n+1}$ 
a lift of $\gamma$.
By construction, this outer action is compatible with the outer action of $\on{Mod}_{g,n}$ on
$\pi_1(C^\circ)$ of \autoref{definition:mcg-finite}.

To show $\rho|_{\pi_1(C^\circ)}$ is MCG-finite,
it therefore suffices to show the conjugacy class of $\rho|_{\pi_1(C^\circ)}$ is stable under
the image $\Gamma$ of $$\pi_1(\mathscr{C}^\circ)\twoheadrightarrow \pi_1(\mathscr{M})\to \pi_1(\mathscr{M}_{g,n})\simeq \on{PMod}_{g,n}.$$
This image has finite index by \autoref{lemma:finite-index-from-dominant-map}. 
The conjugacy class of $\rho|_{\pi_1(C^\circ)}$ is stable under this image because, 
if $\widetilde{\gamma} \in \pi_1(\mathscr{C}^\circ)$ is any lift of  
$\gamma \in \Gamma$, then $\rho^{\widetilde{\gamma}}$ is conjugate to $\rho$ by $\rho(\widetilde{\gamma})$, as in the proof of \autoref{proposition:MCG-finite-index-rep}.
\end{proof}
\subsection{Projective local systems associated to MCG-finite representations}
\label{subsection:mcg-finite-construction}
\source{canonical-representations-surface-groups}

We next give a construction lifting an irreducible MCG-finite representation
of $\pi_1(\Sigma_{g,n})$ to a projective representation of a finite index subgroup of $\Mod_{g,n+1}$.

\begin{notation}
\source{canonical-representations-surface-groups}
\notready
	\label{notation:finite-orbit}
\source{canonical-representations-surface-groups}
	Suppose $\rho$ is an irreducible MCG-finite representation of $\pi_1(\Sigma_{g,n},x)$.
	Let $\Gamma \subset \on{PMod}_{g,n}$ denote the finite index stabilizer
	of the conjugacy class of $\rho$ and let $\widetilde{\Gamma} \subset
	\on{PMod}_{g,n+1}$ denote the preimage of $\Gamma$ under the surjection
	$\on{PMod}_{g,n+1} \to \on{PMod}_{g,n}$ coming from the Birman exact
	sequence.
%	\cite[Theorem 4.6]{farbM:a-primer}.
\end{notation}
	
\begin{lemma}
\source{canonical-representations-surface-groups}
\notready
	\label{lemma:mcg-rep-construction}
\source{canonical-representations-surface-groups}
	In the setup of \autoref{notation:finite-orbit}, there exists a unique representation $\widetilde{\rho}:
	\widetilde{\Gamma} \to \on{PGL}_r(\mathbb C)$ so that
	\begin{equation}
		\label{equation:mcg-rho-and-tilde-rho-compatibility}
\source{canonical-representations-surface-groups}
		\begin{tikzcd} 
			\pi_1(\Sigma_{g,n},x)  \ar[r, "\rho"] \ar {d} &
			 \on{GL}_r(\mathbb C) \ar {d} \\
			 \widetilde{\Gamma}  \ar [r, "\widetilde \rho"] & \on{PGL}_r(\mathbb C)
	\end{tikzcd}\end{equation}
	commutes,
	where $\pi_1(\Sigma_{g,n}, x) \subset\widetilde{\Gamma}\subset \on{PMod}_{g,n+1}$ is the inclusion
	of the point-pushing subgroup from the Birman exact sequence.
%	\cite[Theorem 4.6]{farbM:a-primer}.
\end{lemma}
\begin{proof}
	First, we construct the representation $\widetilde{\rho}$. 
	By identifying 
	$\on{PMod}_{g,n+1} \subset \on{Aut}(\pi_1(\Sigma_{g,n},x))$, 
	we obtain an action of 
	$\on{PMod}_{g,n+1}$ on $r$-dimensional representations of 
	$\pi_1(\Sigma_{g,n},x)$.
	Since $\widetilde{\Gamma}$ preserves $\rho$ up to
	conjugacy, for every $\gamma \in \widetilde{\Gamma}$, there is some matrix $M_\gamma \in \on{GL}_r(\mathbb C)$
	so that $M_\gamma \rho M_\gamma^{-1} = \gamma(\rho)$.
	Since $\rho$ is irreducible, $M_\gamma$ is unique up to scaling by
	Schur's lemma.
We let $\overline{M}_\gamma$ denote the image of $M_\gamma$ in
	$\on{PGL}_r(\mathbb C)$; $\overline{M}_\gamma$ is well-defined by the previous sentence.
	Define $\widetilde{\rho}$ by  
	\begin{align*}
		\widetilde{\rho}:  \widetilde{\Gamma} &\rightarrow \on{PGL}_r(\mathbb C) \\
		\gamma &\mapsto \overline{M}_\gamma.
	\end{align*}

	Uniqueness of 
	$\overline{M}_\gamma \in \on{PGL}_r(\mathbb C)$
	implies that $\widetilde{\rho}$ is a representation, as $\overline{M}_\gamma\overline{M}_{\gamma'}=\overline{M}_{\gamma\gamma'}$. 

	Commutativity of 	
	\eqref{equation:mcg-rho-and-tilde-rho-compatibility}
	follows from the definition of the inclusion $\iota: \pi_1(\Sigma_{g,n}, x) \hookrightarrow
	\on{PMod}_{g,n+1}\subset \on{Aut}(\pi_1(\Sigma_{g,n}, x))$ as the
	point-pushing subgroup, realizing the set of inner automorphisms.
  Indeed, for $\eta\in \pi_1(\Sigma_{g,n}, x)$, $\iota$ sends
	$\eta$ to the automorphism of $\pi_1(\Sigma_{g,n}, x)$ given by $\beta\mapsto \eta \beta
	\eta^{-1}.$
	Hence we may take $M_{\iota(\eta)}=\rho(\eta)$, implying 
	\eqref{equation:mcg-rho-and-tilde-rho-compatibility} commutes.
	
	Finally, uniqueness of $\widetilde{\rho}$ follows from commutativity of \eqref{equation:mcg-rho-and-tilde-rho-compatibility} and uniqueness of $\overline{M}_\gamma$. Indeed, for any $\gamma\in \widetilde{\Gamma}$ and for all $\eta\in \pi_1(\Sigma_{g,n}, x)$, we must have $$\widetilde{\rho}(\gamma\eta\gamma^{-1})=\widetilde{\rho}(\gamma)\widetilde{\rho}(\eta)\widetilde{\rho}(\gamma)^{-1},$$ and hence $\widetilde{\rho}(\gamma)$ must equal $\overline{M}_\gamma$.
\end{proof}
See e.g.~the proof of \cite[Theorem 4]{simpson1992higgs} for a similar argument.

We will also need a geometric variant of \autoref{lemma:mcg-rep-construction}.
We give the geometric rephrasing of \autoref{lemma:mcg-rep-construction},
implicitly using \autoref{lemma:mgn-fundamental-group}. Below $\mathscr{M}$ is
the finite \'etale cover of $\mathscr{M}_{g,n}$ corresponding to a finite-index subgroup of $\Gamma\subset \on{PMod}_{g,n}$, chosen so that $\mathscr{M}$ is a scheme (as opposed to a Deligne-Mumford stack).

\begin{lemma}
\source{canonical-representations-surface-groups}
\notready
	\label{lemma:geometric-mcg-rep-construction}
\source{canonical-representations-surface-groups}
	Let
	$\rho: \pi_1(\Sigma_{g,n},x) \to \on{GL}_r(\mathbb C)$ be an irreducible MCG-finite representation. 
	There exists a scheme $\mathscr M$ with a finite \'etale map $\mathscr M \to \mathscr
	M_{g,n}$, with associated family of $n$-times punctured curves $\pi^\circ: \mathscr{C}^\circ\to \mathscr{M}$, a point $c\in \mathscr{C}^\circ$ with $m:=\pi^\circ(c)$,
	and a representation $\widetilde{\rho}: \pi_1(\mathscr C^\circ, c) \to
	\on{PGL}_r(\mathbb C)$ so that the following holds. Upon identifying
	$\pi_1(\Sigma_{g,n},x)  \simeq \pi_1(\mathscr C^\circ_m,c)$, the diagram
	\begin{equation}
		\label{equation:geometric-rho-and-tilde-rho}
\source{canonical-representations-surface-groups}
		\begin{tikzcd} 
			\pi_1(\mathscr C^\circ_m,c)  \ar[r, "\rho"] \ar {d} &
			 \on{GL}_r(\mathbb C) \ar {d} \\
			 \pi_1(\mathscr C^\circ, c) \ar[r, "\widetilde\rho"]& \on{PGL}_r(\mathbb C)
	\end{tikzcd}\end{equation}
	commutes.
\end{lemma}
%\begin{proof}
%	Since $\Gamma \subset \on{Mod}_{g,n}$ is a finite index subgroup,
%	and $\on{Mod}_{g,n} \simeq\pi_1(\mathscr M_{g,n})$
%	\autoref{lemma:mgn-fundamental-group}
%	we obtain a corresponding finite \'etale cover $\mathscr M$ of $\mathscr M_{g,n}$.
%	It may be that $\mathscr M$ fails to be a scheme, but we can make it a
%	scheme by passing to a further finite \'etale cover. (For example, the
%	cover obtained by specifying a basis for the $5$-torsion of the Jacobian
%	of the universal curve is a scheme).
%	After possibly passing to this further cover, and taking
%	$\widetilde{\Gamma}$ to be the preimage of $\Gamma$ in
%	$\on{Mod}_{g,n+1}$ as in \autoref{lemma:mcg-rep-construction}
%	We obtain an isomorphism $\widetilde{\Gamma} \simeq \pi_1(\mathscr
%	C^\circ,c)$.
%	Under these identifications,
%	the existence of $\widetilde{\rho}$ and commutativity of
%	\eqref{equation:geometric-rho-and-tilde-rho} 
%	follow from \autoref{lemma:mcg-rep-construction}.
%\end{proof}



\subsection{Lifting Projective representations}
\label{subsection:lifting-projective}
\source{canonical-representations-surface-groups}

The main goal of this section is to prove \autoref{proposition:gl-lift},
which says that, \'etale locally on the base $\mathscr{M}$, we can lift the representation $\widetilde{\rho}$ into
$\on{PGL}_r(\mathbb C)$
constructed in \autoref{lemma:geometric-mcg-rep-construction}
to a representation $\rho'$ into $\on{GL}_r(\mathbb C)$.

Because the obstruction to lifting is related to $\mathbb C^\times =
\ker(\on{GL}_r(\mathbb C) \to
\on{PGL}_r(\mathbb C))$, we will use the following classification when
$r = 1$.

\begin{proposition}[\protect{~\cite[Lemma 3.2]{biswas2017surface}}]
\source{canonical-representations-surface-groups}
\notready
	\label{proposition:1-dim-finite-image}
\source{canonical-representations-surface-groups}
	Let $g\geq 1, n\geq 0$. Then any MCG-finite representation $$\rho: \pi_1(\Sigma_{g,n})\to \mathbb{C}^\times$$ has finite image.
\end{proposition}
For the reader's benefit, we recall the idea of the proof of
\autoref{proposition:1-dim-finite-image}.
The idea is to show
that if any standard generator of $\pi_1(\Sigma_{g,n})$ corresponding to a loop
$\gamma$ in $\Sigma_{g,n}$ has infinite order under $\rho$, then acting on $\rho$ by powers of the
Dehn twist about $\gamma$ gives an infinite collection of distinct
representations.

In order to lift 
$\on{PGL}_r(\mathbb C)$-reps to
$\on{GL}_r(\mathbb C)$-reps, 
we will need to know 
that the image of $\rho$ has finite
intersection with the center of $\on{GL}_r(\mathbb C)$.
%as we now verify using \autoref{proposition:1-dim-finite-image}.
\begin{lemma}
\source{canonical-representations-surface-groups}
\notready
	\label{lemma:finite-central-part}
\source{canonical-representations-surface-groups}
	Let $g\geq 1$.
	Suppose $\rho: \pi_1(\Sigma_{g,n}) \to \on{GL}_r(\mathbb C)$ is an MCG-finite representation.
	Then $\rho$ has finite determinant.
	Moreover, $\rho$ factors through a subgroup $G$ with $\on{SL}_r(\mathbb C) \subset G
	\subset \on{GL}_r(\mathbb C)$ and $\mu := G \cap \ker(\on{GL}_r(\mathbb
	C) \to \on{PGL}_r(\mathbb C))$ a finite group.
\end{lemma}
\begin{proof}
	By \autoref{proposition:1-dim-finite-image}, the composition
	$\pi_1(\Sigma_{g,n}) \xrightarrow{\rho} \on{GL}_r(\mathbb C) \xrightarrow{\on{det}}
	\mathbb C^\times$ has finite image $H$.
	We can take $G := \det^{-1}(H) \subset \on{GL}_r(\mathbb{C})$,
	which indeed has finite intersection with the center of
	$\on{GL}_r(\mathbb C)$, and hence finite intersection with
	$\ker(\on{GL}_r(\mathbb C) \to \on{PGL}_r(\mathbb C))$
\end{proof}

In order to lift representations from
$\on{PGL}_r(\mathbb C)$ to
$\on{GL}_r(\mathbb C)$,
we will need to kill certain cohomological obstructions using dominant \'etale maps.
\begin{lemma}
\source{canonical-representations-surface-groups}
\notready
	\label{lemma:vanish-on-etale}
\source{canonical-representations-surface-groups}
	Let $i > 0$ and $\mathscr{M}$ be a complex variety. Suppose $\mu$ is a finite abelian group and
	$\sigma \in H^i(\pi_1(\mathscr M, m), \mu)$
	a cohomology class.
	Then there is a dominant \'etale map $\mathscr M' \to \mathscr M$
	so that $\sigma|_{\mathscr M'} =0 \in H^i(\pi_1(\mathscr M'), \mu)$.
\end{lemma}
\begin{proof}
%Alternate proof:
%	Let us use $X$ for a topological space representing $K(\pi_1(\mathscr
%	M,m), 1)$.
%	We obtain an induced map $f: \mathscr M \to X$
%	and a corresponding spectral sequence with low degree terms
%\begin{align*}
%	0 \to H^1(X, f_* \mu) \to H^1(\mathscr M, \mu) \to H^0(X, R^1 f_* \mu)
%	\to H^2(X, f_* \mu) \to H^2(\mathscr M, \mu).
%\end{align*}
%
%We claim $R^1 f_* \mu = 0$.
%Indeed, if $\widetilde{X}$ is the universal cover of $X$ and $\widetilde{\mathscr M}$ is the
%universal cover of $M$,
%we find $\widetilde{\mathscr M} \to \widetilde{X}$ is the pullback of $\mathscr M \to
%X$, and so the stalk of $R^1 f_* \mu$ at any point of $\widetilde{X}$ is identified with
%$H^1(\widetilde{\mathscr M}, \mu) =
%\Hom(\pi_1(\widetilde{\mathscr M}), \mu) = 0$.  

First, we claim that for any class
$\alpha \in H^i(\mathscr M, \mu)$,
there is a dominant \'etale map $\mathscr N \to \mathscr M$ so that
$\alpha|_{\mathscr N} = 0$.
Since $\mu$ is finite, we can use the comparison between singular
and \'etale cohomology to reduce to showing that for any
$\alpha \in H^i_{\mathrm{\acute et}}(\mathscr M, \mu)$ there is a dominant \'etale map
$\mathscr N \to \mathscr M$ for which $\alpha|_{\mathscr N}= 0$.
Choose an injective resolution $I^\bullet$ of $\mu$ and represent
$\alpha$ by some $\bar\alpha \in \ker(I^i(\mathscr M) \to I^{i+1}(\mathscr M))$.
Because $I^\bullet$ is exact, there exists for each $x\in \mathscr{M}$ an \'etale neighborhood $U_x$ of $x$  
so that $\alpha|_{U_x}$ is in the image of $I^{i-1}(U_x)\to I^{i}(U_x)$. Taking $x$ to be the generic point of $\mathscr{M}$ gives the claim.

Now, let $\alpha$ as above denote the image of $\sigma$ under the natural map
$H^i(\pi_1(\mathscr M, m), \mu) \to H^i(\mathscr M, \mu),$
and take $\mathscr N$ as above so that $\alpha|_\mathscr N = 0$.
By
\cite[Expos\'e XI, 4.6]{SGA4}
we can find a Zariski open $\mathscr M'
\subset \mathscr N$ which is a 
$K(\pi_1(\mathscr M', m'), 1)$ Eilenberg-Maclane space.
This implies the natural map $H^i(\pi_1(\mathscr M', m'), \mu) \to H^i(\mathscr M', \mu)$
is an isomorphism.
Therefore, since $\alpha|_{\mathscr N} = 0$, we also have $\alpha|_{\mathscr M'}
= 0$, and hence $\sigma|_{\pi_1(\mathscr M')} = 0$, as desired.
\end{proof}

We now combine the above results to show that the projective representation $\widetilde{\rho}$ constructed in \autoref{lemma:geometric-mcg-rep-construction}
can be lifted to a linear representation.

\begin{proposition}
\source{canonical-representations-surface-groups}
\notready
	\label{proposition:gl-lift}
\source{canonical-representations-surface-groups}
	Suppose $\mathscr C \to \mathscr M$ is a versal family of $n$-pointed
	curves of genus $g\geq 1$, with associated family of punctured curves 
	$\pi^\circ: \mathscr{C}^\circ\to \mathscr{M}$. Let $c\in \mathscr{C}^\circ$ be a point and $m=\pi^\circ(c)$.
	Suppose we have representations $\rho: \pi_1(\mathscr C_m^\circ, c) \to
	\on{GL}_r(\mathbb C)$
	and $\widetilde{\rho}: \pi_1(\mathscr C^\circ, c) \to \on{PGL}_r(\mathbb
	C)$ so that the diagram
	\begin{equation}
		\label{equation:rho-and-tilde-rho-compatibility}
\source{canonical-representations-surface-groups}
		\begin{tikzcd} 
			 \pi_1(\mathscr C_m^\circ, c) \ar[r, "\rho"] \ar {d} &
			 \on{GL}_r(\mathbb C) \ar {d} \\
			 \pi_1(\mathscr C^\circ, c) \ar[r, "\widetilde\rho"] & \on{PGL}_r(\mathbb C)
	\end{tikzcd}\end{equation}
	commutes.
	Moreover, assume $\rho$ is MCG-finite.
	Then, there exists 
	\begin{enumerate}
		\item a dominant \'etale map $\mathscr M' \to \mathscr M$,
		\item corresponding relative curve $\mathscr C' := \mathscr M'
	\times_{\mathscr M} \mathscr C$ with associated family of punctured curves ${\mathscr{C}'}^\circ$,
	and 
\item a representation $\rho' : \pi_1(\mathscr C'^\circ, c') \to
	\on{GL}_r(\mathbb C)$ for some basepoint $c' \in \mathscr C'^\circ$ lying over $c$
	\end{enumerate}
so that, 
	upon choosing some $m'$ over $m$ with $\mathscr C_{m'} \simeq \mathscr C_m$,
	\begin{equation}
		\label{equation:cover-rep-compatibility}
\source{canonical-representations-surface-groups}
		\begin{tikzcd} 
			\pi_1(\mathscr C_{m'}^\circ, c') \ar{r} \ar[bend left]{rr}{\rho} &
			\pi_1(\mathscr C'^\circ, c') \ar {r}{\rho'} \ar {d} &
			 \on{GL}_r(\mathbb C) \ar {d} \\
			 & \pi_1(\mathscr C^\circ, c) \ar
			 {r}{\widetilde{\rho}} & \on{PGL}_r(\mathbb C)
	\end{tikzcd}\end{equation}
	commutes.	
	Moreover, $\on{im} \rho'$ is contained a subgroup $G$ with $\on{SL}_r(\mathbb
	C) \subset G \subset
	\on{GL}_r(\mathbb C)$, such that $\on{SL}_r(\mathbb C) \subset G$ has
	finite index.
\end{proposition}
\begin{proof}
	Taking $G$ and $\mu$ as in \autoref{lemma:finite-central-part},
	there is a commuting diagram
\begin{equation}
	\label{equation:pgl-and-birman}
\source{canonical-representations-surface-groups}
	\begin{tikzcd} 
H^1(\pi_1(\mathscr C^\circ,c), \mu) \ar
		{r}{\lambda} \ar
		{d}{\nu} &  H^1(\pi_1(\mathscr
				C_m^\circ,c),
		\mu) \ar {d}{\xi} \\
		\Hom(\pi_1(\mathscr C^\circ,c), G) \ar
		{r}{\alpha} \ar
		{d}{\beta} & \Hom(\pi_1(\mathscr
				C_m^\circ,c),
			G) \ar {d}{\gamma} \\
		\Hom(\pi_1(\mathscr C^\circ,c), \on{PGL}_r(\mathbb C)) \ar
		{r}{\delta}\ar{d}{\varepsilon}
		&\Hom(\pi_1(\mathscr C_m^\circ,c),
		\on{PGL}_r(\mathbb C))\ar{d}{\zeta} \\
			H^2(\pi_1(\mathscr C^\circ,c), \mu) \ar {r}{\eta}
		& H^2(\pi_1(\mathscr C_m^\circ,c),
		\mu)
\end{tikzcd}\end{equation}
where the columns are exact sequences of sets.
We start with a representation $\widetilde{\rho} \in \on{Hom}(\pi_1(\mathscr C^\circ,c),
\on{PGL}_r(\mathbb C))$ 
and $\rho \in \on{Hom}(\pi_1(\mathscr C_m^\circ, c), G)$.
The commutativity of \eqref{equation:rho-and-tilde-rho-compatibility} can be equivalently
expressed as the statement that $\delta(\widetilde{\rho}) = \gamma(\rho)$. 
We seek to construct $\rho'$ such that \eqref{equation:cover-rep-compatibility} commutes,
meaning that, after replacing $\mathscr M$ with a dominant \'etale $\mathscr M'$ and $\mathscr{C}$ with $\mathscr{C}'$,
there is some $\rho'$ with $\beta(\rho') = \widetilde{\rho}$ and $\alpha(\rho')
= \rho$.

First, we argue we can pass to some $\mathscr M'$ so that $\widetilde{\rho}$
lies in the image of $\beta$. 
It is equivalent to arrange that $\varepsilon(\widetilde{\rho}) = 0$.
As a first step, we claim $\eta(\varepsilon(\widetilde{\rho})) = 0$.
This holds because $\eta(\varepsilon(\widetilde{\rho})) =
\zeta(\delta(\widetilde{\rho}))$, and
$\delta(\widetilde{\rho}) = \gamma(\rho)$ lies in the image of $\gamma$.
We therefore have that $\varepsilon(\widetilde{\rho}) \in \ker \eta$. 
The Hochschild-Serre spectral sequence associated to the short exact sequence  $$1\to \pi_1(\mathscr{C}_m^\circ)\to \pi_1(\mathscr{C}^\circ)\to \pi_1(\mathscr{M})\to 1$$ of
\autoref{lemma:SES-pi1-versal-curves} yields a short exact sequence $$0\to H^{2,0}\to \ker \eta \overset{\iota}{\to} H^{1,1}\to 0,$$ where $H^{i,j}$ is a subquotient of $H^i(\pi_1(\mathscr M, m), H^j(\pi_1(\mathscr C^\circ_m,c),
\mu))$.

By passing to a finite \'etale cover $\mathscr{M}_1$ of $\mathscr{M}$, we can assume $\pi_1(\mathscr{M}, m)$ acts trivially on the finite group $H^1(\pi_1(\mathscr{C}_m^\circ, c), \mu)$.  Now as $H^i(\pi_1(\mathscr M, m), H^j(\pi_1(\mathscr C^\circ_m,c),
\mu))$ is finite, we may by \autoref{lemma:vanish-on-etale} pass to a further
dominant \'etale scheme over $\mathscr{M}$ killing every element of $H^{i,j}$,
$(i,j)=(2,0) \text{ or } (1,1)$, and hence every element of $\ker \eta$. Thus
after passing to any component $\mathscr{M}_2$ of this dominant \'etale scheme
over $\mathscr M_2$, we have $\epsilon(\widetilde\rho)=0$. Let $\mathscr{C}_2^\circ$ be the pullback of $\mathscr{C}^\circ$ to $\mathscr{M}_2$. 

Let $\overline{\rho} \in H^1(\pi_1(\mathscr C_2^\circ, c), G)$
be an element with $\beta(\overline{\rho}) = \widetilde{\rho}$; such a $\overline{\rho}$ exists as we have arranged $\epsilon(\widetilde{\rho})=0$. 
We wish to arrange $\alpha(\overline{\rho}) =\rho$. 
This may not be the case, but it is enough to show that after passing to a further dominant
\'etale $\mathscr M'$,
that we can modify $\overline{\rho}$ by an element of $\on{im}\nu$ so that this
does hold.
More precisely, note that $\alpha(\overline{\rho})$ differs from $\rho$ by an
element of $\sigma\in H^1(\pi_1(\mathscr{C}^\circ_m, c), \mu)$, as by
construction $\gamma(\alpha(\overline{\rho}))=\gamma(\rho)$. It suffices to show
that, after passing to a dominant \'etale scheme over $\mathscr{M}_2$, $\sigma$ is in the image of $\lambda$. After passing to a finite \'etale cover of $\mathscr{M}_2$, we may assume that $H^1(\pi_1(\mathscr{C}^\circ_m, c), \mu)$ is fixed by $\pi_1(\mathscr{M}, m)$.
The Hochschild-Serre spectral sequence gives a map
\begin{align*}
	H^0(\pi_1(\mathscr M, m), H^1(\pi_1({\mathscr C}_{m}^\circ,c), \mu))
\xrightarrow{\chi}
H^2(\pi_1(\mathscr M, m), H^0(\pi_1({\mathscr C}^\circ_{m}, c), \mu))
\end{align*}
such that $\chi(\sigma) = 0$ if an only if $\sigma \in \on{im}\lambda$.
Since $H^0(\pi_1({\mathscr C}^\circ_{m}, c), \mu)=\mu$
is finite, we can apply
\autoref{lemma:vanish-on-etale} so as to assume, after replacing $\mathscr M$ by
a dominant \'etale $\mathscr M'$
that $\chi(\sigma) = 0$.
This implies that $\sigma \in \on{im}\lambda$, as desired.
\end{proof}

We can now give a partial converse to \autoref{proposition:MCG-finite-family-of-curves}.
\begin{corollary}
\source{canonical-representations-surface-groups}
\notready\label{corollary:converse-to-MCG-rep}
\source{canonical-representations-surface-groups}
Let $\rho:\pi_1(\Sigma_{g,n})\to \on{GL}_r(\mathbb{C})$ be an irreducible MCG-finite
representation, with $g\geq 1$. 
There is a punctured versal family of curves
$\mathscr{C}^\circ\to \mathscr{M}$, and a representation $\rho'$ of
$\pi_1(\mathscr{C}^\circ)$ whose determinant has finite order, with the following property: For $C^\circ$ a fiber of
$\pi^\circ$, there is an identification $\pi_1(\Sigma_{g,n})\simeq
\pi_1(C^\circ)$ such that, under this identification, $\rho'|_{\pi_1(C^\circ)}=\rho$.
\end{corollary}
\begin{proof}
	This follows by combining \autoref{lemma:geometric-mcg-rep-construction} and \autoref{proposition:gl-lift}.
\end{proof}

\begin{remark}
\source{canonical-representations-surface-groups}
\notready
	\label{remark:cousin-heu}
\source{canonical-representations-surface-groups}
	Combining \autoref{corollary:converse-to-MCG-rep} with
	\autoref{proposition:MCG-finite-family-of-curves} gives a proof of
	\cite[Theorem A]{CH:isomonodromic}
	in the case that $\rho$ is semisimple and the flat vector bundle $(E, \nabla_0)$
	of \cite[Theorem A]{CH:isomonodromic} 
	is taken to be the Deligne canonical extension of
	$(\rho|_{C^\circ}\otimes \mathscr{O}, \on{id}\otimes d)$ to $C$.
\end{remark}



\subsection{Representations that are unitary on fibers}
\label{subsection:unitary-on-fibers}
\source{canonical-representations-surface-groups}

The main goal of this section is to verify
\autoref{lemma:unitary-direct-sum}, which allows us to write local systems on a family of punctured curves $\mathscr{C}^\circ\to \mathscr{M}$ in terms of unitary local systems and local systems pulled back from $\mathscr{M}$, after passing to a dominant \'etale
map.

The following criterion for unitarity and finiteness will be useful when analyzing the representations produced by \autoref{lemma:mcg-rep-construction} and \autoref{proposition:gl-lift}.
\begin{lemma}
\source{canonical-representations-surface-groups}
\notready
	\label{lemma:finite-rho-implies-finite-lift}
\source{canonical-representations-surface-groups}
Let $G$ be a group and $H\trianglelefteq G$ a normal subgroup. Let $$\rho: G\to \on{GL}_r(\mathbb{C})$$ be a representation such that $\det\rho$ has finite order, and suppose $\rho|_H$ is irreducible.
\begin{enumerate}
\item If $\rho|_H$ has finite image, then $\rho$ has finite image.
\item if $\rho|_H$ is unitary, then $\rho$ is unitary.
\end{enumerate}
\end{lemma}
\begin{proof}
	We first prove $(1)$.
		As $\det\rho$ has finite order, it suffices to show that the
	projectivization $$\mathbb{P}\rho: G\overset{\rho}{\longrightarrow}
	\on{GL}_r(\mathbb{C})\to \on{PGL}_r(\mathbb{C})$$ has finite image. Let
$t=\#\on{im}(\rho|_H)$. For each $g\in G$, $\mathbb{P}\rho(g)$ is the unique
(by Schur's lemma) element of $\on{PGL}_r(\mathbb{C})$ such that
$$\mathbb{P}\rho(g)\mathbb{P}\rho(h)\mathbb{P}\rho(g)^{-1}=\mathbb{P}\rho(ghg^{-1})$$ for all $h\in H$.
Since $\rho(g)$ acts by conjugation on the order $t$ set $\on{im}(\rho|_H)$, its action 
has order dividing $t!$, so $\rho(h)=\rho(g^{t!})\rho(h)\rho(g^{-t!})$. Hence we have $\mathbb{P}\rho(g^{t!})=\on{id}$ by
uniqueness. Thus the image of $\mathbb{P}\rho$ has exponent dividing $t!$. But a linear group with finite exponent is finite \cite{burnside1905criteria}.

We now turn to the unitary case and verify $(2)$.
Let $h$ be a positive-definite Hermitian inner product preserved by $\rho|_H$. 
We first check $\on{im}\rho$ lies in the general unitary group, i.e.~it preserves $h$ up to scaling.
	Let $V$ be the underlying vector space of our
	representation $\rho$. The invariant inner product $h: V \times V \to \mathbb C$ corresponds to a $\mathbb{C}$-linear isomorphism of $H$-representations $\phi: V \to \overline{V}^\vee$.
	The linear map $\phi$, and hence $h$, is unique up to scaling by Schur's lemma. 
	For any $g\in G$, $$h^g: (v, w)\mapsto h(\rho(g)v, \rho(g)w)$$ is
	another $H$-invariant Hermitian form, and hence $h^g=c\cdot h$ for some
	$c\in\mathbb{C}^\times.$ 
	Equivalently, $\on{im}(\rho)\subset GU(h)$. 
	Since $\rho$ has finite determinant, we moreover obtain
	$\on{im}(\rho)\subset U(h)$.\qedhere
\end{proof}
We conclude the section by giving a convenient description of local systems
$\mathbb V$ on a family of punctured curves $\pi^\circ: \mathscr C^\circ \to
\mathscr M$, with unitary fibral monodromy. It may be useful for the reader to consider the case $A=\mathbb{C}$.

As in \autoref{notation:versal-family}, let $\pi: \mathscr{C}\to \mathscr{M}$ be a versal family of $n$-pointed curves, and let $\mathscr{C}^\circ\to \mathscr{M}$ be the associated family of punctured curves. Let $m\in\mathscr{M}$ be a point.

\begin{lemma}
\source{canonical-representations-surface-groups}
\notready
	\label{lemma:unitary-direct-sum}
\source{canonical-representations-surface-groups}
	Let $A$ be an Artin local $\mathbb{C}$-algebra with residue field $\mathbb{C}$.
	Suppose $\mathbb V$ is a local system of free $A$-modules on $\mathscr C^\circ$
	such that $\mathbb V|_{\mathscr C^\circ_m}$ is a constant deformation of a unitary local system, i.e.~there exists a unitary $\mathbb{C}$-local system $\mathbb{V}_0$ on $\mathscr C^\circ_m$ such that $\mathbb V|_{\mathscr C^\circ_m}\simeq \mathbb{V}_0\otimes_\mathbb{C} A$.  
	
	There is a dominant \'etale map $\mathscr M' \to \mathscr M$, with $\mathscr{C}'=\mathscr{M}'\times_\mathscr{M}\mathscr{C}$ and ${\pi'}^\circ: {\mathscr{C}'}^\circ\to \mathscr{M}'$  the associated family of punctured curves, over
	which 
	\begin{align*}
	\mathbb V|_{ {\mathscr{C}'}^\circ} \simeq \oplus_{i=1}^s \mathbb U_i
\otimes ({\pi'}^\circ)^* \mathbb
W_i,
	\end{align*}
where the $\mathbb W_i$ are locally constant sheaves of free $A$-modules on $\mathscr M'$ and
	the $\mathbb U_i$ are unitary local systems on ${\mathscr{C}'}^\circ$.
	Moreover, for $C'^\circ$ a fiber of $\mathscr{C}'^\circ \to
	\mathscr{M}'$, each $\mathbb{U}_i|_{C'^\circ}$ is irreducible, the
	$\mathbb{U}_i|_{C'^\circ}$ are pairwise non-isomorphic,
	and $\mathbb W_i \simeq {\pi'}^\circ_*\Hom(\mathbb U_i, \mathbb V|_{ {\mathscr{C}'}^\circ} )$.
\end{lemma}
\begin{proof}
Since $\mathbb V_0$ is unitary, we can express it as
as a sum of irreducible unitary local systems, 
$\mathbb V_0 \simeq \oplus_{i=1}^s \mathbb
S_i^{\oplus n_i}$, with the $\mathbb{S}_i$ pairwise non-isomorphic.
Let $\rho:\pi_1(\mathscr C_m^\circ) \to \on{GL}_r(\mathbb C)$ be the monodromy representation associated to $\mathbb{V}_0$, and let
$\rho_i$ be the (irreducible) representation associated to $\mathbb{S}_i$.
By \autoref{proposition:MCG-finite-family-of-curves}, $\rho$ is MCG-finite. Hence each
$\rho_i$ is MCG-finite by \autoref{proposition:basic-properties}.

By repeatedly applying \autoref{corollary:converse-to-MCG-rep},
there is a dominant \'etale map $\mathscr M' \to \mathscr M$ and representations
$\rho'_i: \pi_1({\mathscr {C}'}^\circ) \to \on{GL}_r(\mathbb C)$ with finite determinant so that for any $m' \in \mathscr
M'$, $\rho'_i$ restricts to a representation
$\pi_1({\mathscr{C}'}^\circ_{m'}) \to \on{GL}_r(\mathbb C)$
identified with $\rho_i$.
Let $\mathbb U_i$ denote the local system on ${\mathscr{C}'}^\circ$ corresponding to
$\rho_i'$.
Each $\mathbb U_i$ is unitary by \autoref{lemma:finite-rho-implies-finite-lift}.
The $\mathbb{U}_i$ 
are irreducible and pairwise non-isomorphic
because the same holds for the $\rho_i$.

Let $\mathbb W_i = \pi'^\circ_*\Hom(\mathbb U_i, \mathbb V|_{\mathscr
C'^\circ})$, as in the statement. The $\mathbb{W}_i$ are locally constant sheaves of free $A$-modules by the hypothesis that $\mathbb V|_{\mathscr C^\circ_m}\simeq \mathbb{V}_0\otimes_\mathbb{C} A$. There is a natural evaluation map
\begin{align*}
	\psi: \oplus_{i=1}^s (\pi'^\circ)^* \mathbb W_i \otimes \mathbb U_i \to \mathbb
	V|_{{\mathscr{C}'}^\circ}.
\end{align*}
It only remains to show $\psi$ is an isomorphism.
We do so after restriction to any fiber
of $\mathscr C'^\circ \to \mathscr M'$, where $\psi$ is identified with the
isomorphism
\begin{align}
	\nonumber
	\oplus_{i=1}^s \Hom(\rho_i, \rho\otimes A) \otimes \rho_i  & \rightarrow \rho\otimes A\\
	\sum_{i=1}^s \alpha_1 \otimes v_i & \mapsto \sum_{i=1}^s \alpha_i(v_i).
	\qedhere
\end{align}
%This is an isomorphism because $\rho$ is unitary, hence semisimple.
%Here $\Hom(\rho_i, \rho)$ denotes the vector space of maps from $\rho_i$
%to $\rho$.
%This is indeed an isomorphism as the $\rho_i$ account for 
%all irreducible representations appearing in $\rho$, by definition.
\end{proof}
